% ==========================================
% BAB V RENCANA SELANJUTNYA
% ==========================================
\chapter{RENCANA SELANJUTNYA}
\label{chap:rencana-selanjutnya}
% Pada bab ini dijelaskan secara detail langkah-langkah rencana selanjutnya, hal-hal yang diperlukan atau akan disiapkan, dan risiko dan mitigasinya.

% \begin{enumerate}
% \item	Rencana implementasi, termasuk alat dan bahan yang diperlukan, lingkungan, konfigurasi, biaya, dan sebagainya.
% \item	Desain pengujian dan evaluasi, misalnya metode verifikasi dan validasi.
% \item	Analisis risiko dan mitigasi, misalnya tindakan selanjutnya jika ada yang tidak berjalan sesuai rencana.
% \end{enumerate}
Bab ini menjelaskan rencana kegiatan lanjutan yang akan dilakukan setelah tahap perancangan konsep solusi. Rencana selanjutnya mencakup pengembangan desain sistem secara lebih rinci, implementasi perangkat lunak audit energi, pengujian dan evaluasi sistem, serta penyusunan dokumentasi akhir tugas akhir. Seluruh tahapan dirancang secara sistematis dan berurutan sesuai dengan metodologi pengembangan perangkat lunak yang telah ditetapkan, sehingga diharapkan mampu menghasilkan sistem yang fungsional, teruji, dan sesuai dengan tujuan penelitian. Bab ini juga menjelaskan secara hal-hal yang diperlukan atau akan disiapkan, serta risiko dan mitigasinya.

\section{Pengembangan Desain Sistem Lanjutan}
Pengembangan desain sistem pada tahap ini difokuskan untuk melengkapi dan memperdalam rancangan sistem audit energi yang telah dibahas pada bab sebelumnya. Setelah desain fungsional sistem didefinisikan melalui \textit{use case}, diperlukan perancangan lanjutan untuk memastikan bahwa sistem dapat diimplementasikan secara konsisten dan sesuai dengan kebutuhan nyata di lapangan. Oleh karena itu, pengembangan desain sistem mencakup pengumpulan kebutuhan lanjutan melalui wawancara dan observasi, perancangan desain data sistem, serta perancangan perilaku sistem.

\subsection{Pengumpulan Kebutuhan Lanjutan}
Pengumpulan kebutuhan lanjutan dilakukan untuk memperdalam pemahaman terhadap proses audit energi yang dijalankan di lapangan serta memastikan kesesuaian rancangan sistem dengan praktik nyata. Kegiatan ini mencakup wawancara dan observasi kepada pihak yang memahami dan terlibat langsung dalam audit energi dan pengelolaan sistem kelistrikan gedung, seperti auditor energi, teknisi pemeliharaan, atau pengelola fasilitas gedung. Wawancara difokuskan pada alur kerja audit yang digunakan, jenis data yang dikumpulkan dan dicatat, kendala yang sering dihadapi selama audit, serta kebutuhan informasi pada tahap analisis dan pelaporan. Observasi dilakukan untuk mengamati secara langsung proses pencatatan dan pengolahan data audit yang berjalan saat ini. Informasi yang diperoleh dari tahap ini digunakan untuk memvalidasi kebutuhan sistem, memperkaya skenario penggunaan, serta menjadi dasar yang kuat dalam perancangan desain data dan desain perilaku sistem pada tahap selanjutnya.

\subsection{Desain Data Sistem}
Desain data sistem bertujuan untuk memodelkan data yang diperlukan dalam pelaksanaan audit energi secara terstruktur dan konsisten. Pada tahap ini, data audit direpresentasikan dalam bentuk entitas, atribut, dan relasi yang mencerminkan objek utama dalam audit energi, seperti gedung, proyek audit, ruangan, data konsumsi energi, indikator kinerja energi, temuan audit, serta hasil penilaian dan laporan. Perancangan desain data dilakukan menggunakan \textit{Entity Relationship Diagram} (ERD) untuk menggambarkan keterkaitan antar data dan memastikan bahwa seluruh informasi yang dibutuhkan selama proses audit dapat disimpan, ditelusuri, dan dikelola secara terpusat. Desain data ini menjadi dasar bagi implementasi basis data sistem dan mendukung kebutuhan penyimpanan historis serta analisis audit energi secara berkelanjutan.

\subsection{Desain Perilaku Sistem}
Desain perilaku sistem digunakan untuk menggambarkan bagaimana sistem audit energi merespons interaksi pengguna dan kejadian tertentu selama proses audit. Perancangan dilakukan dengan memodelkan alur interaksi antara aktor dan sistem pada skenario utama audit energi, seperti penginputan data, validasi dan sinkronisasi data, perhitungan indikator kinerja energi, serta penyajian hasil audit. Diagram perilaku, seperti sequence diagram, digunakan untuk menunjukkan urutan pesan dan proses yang terjadi secara dinamis, termasuk proses otomatis yang dijalankan oleh sistem tanpa interaksi langsung dari pengguna. Desain perilaku ini membantu memastikan bahwa alur proses yang diimplementasikan selaras dengan kebutuhan fungsional serta mendukung otomatisasi proses audit energi secara konsisten dan terkontrol.

\section{Rencana Implementasi Sistem}
Rencana implementasi sistem dilakukan berdasarkan tahapan metodologi Waterfall yang telah dijelaskan pada Bab \ref{chap:pendahuluan}. Pada tahap ini, rancangan desain fungsional, desain data, dan desain perilaku sistem yang telah disusun akan diterjemahkan ke dalam bentuk perangkat lunak yang dapat dijalankan. Implementasi mencakup pengembangan modul input dan pengelolaan data audit, pemrosesan dan perhitungan indikator kinerja energi, penyimpanan data terpusat, serta penyajian hasil audit dalam bentuk \textit{dashboard} dan laporan otomatis. Proses implementasi direncanakan dilakukan secara bertahap dan terstruktur untuk memastikan setiap fungsi berjalan sesuai dengan kebutuhan fungsional yang telah ditetapkan.

\subsection{Lingkungan dan Alat Pengembangan}
Lingkungan dan alat pengembangan disusun untuk mendukung rencana implementasi sistem audit energi yang telah dirumuskan pada subbab sebelumnya. Pemilihan alat dan teknologi difokuskan pada kemudahan pengembangan, dukungan terhadap sistem berbasis web, serta kemampuan penyimpanan data terpusat dengan dukungan sinkronisasi daring. Lingkungan pengembangan ini dirancang agar sesuai dengan skala penelitian dan kebutuhan fungsional sistem, sehingga proses implementasi dapat dilakukan secara efisien dan selaras dengan rancangan desain sistem yang telah ditetapkan.

\input table/tabeltools.tex

\section{Rencana Pengujian dan Evaluasi}
Rencana pengujian dan evaluasi sistem dilakukan untuk memastikan bahwa perangkat lunak audit energi yang dikembangkan telah memenuhi kebutuhan fungsional dan nonfungsional yang telah ditetapkan. Pengujian fungsional direncanakan dengan memverifikasi setiap fungsi utama sistem, seperti impor dan penginputan data, perhitungan indikator kinerja energi, visualisasi hasil audit, serta pembuatan laporan otomatis, berdasarkan kebutuhan fungsional yang telah dirumuskan. Pengujian nonfungsional difokuskan pada aspek kegunaan, konsistensi hasil perhitungan, dan keandalan sistem dalam menangani skenario penggunaan seperti input data secara luring dan sinkronisasi data saat koneksi tersedia, berdasarkan kebutuhan nonfungsional yang telah dirumuskan. Evaluasi hasil sistem dilakukan dengan membandingkan keluaran perangkat lunak terhadap perhitungan manual dan mengecek kepatuhan standar audit energi pada studi kasus satu gedung, guna menilai efektivitas sistem dalam mendukung proses audit energi.

\section{Rencana Waktu Pelaksanaan}
Rencana waktu pelaksanaan tugas akhir ini disusun untuk memberikan gambaran tahapan pekerjaan secara sistematis dan terukur. Perencanaan waktu disajikan menggunakan pendekatan Gantt Chart yang menunjukkan keterkaitan antar tahapan serta alokasi waktu pelaksanaan setiap kegiatan. Setiap kegiatan yang direncanakan sudah selaras dengan tahapan pengembangan perangkat lunak Waterfall. Rincian rencana waktu pelaksanaan ditunjukkan pada Tabel~\ref{tbl:timeline}, yang mencakup seluruh tahapan mulai dari studi awal hingga penyusunan dokumentasi akhir tugas akhir.

\input table/tabeltimeline.tex

Berdasarkan Tabel~\ref{tbl:timeline}, tahapan awal meliputi studi awal dan perumusan topik pada bulan pertama (B1) serta studi literatur pada bulan pertama dan kedua (B1-B2) telah terlaksanakan. Tahapan selanjutnya mencakup analisis kondisi dan kebutuhan sistem serta desain sistem (B2–B3) yang masih perlu untuk diperdalam pemahaman kebutuhan penggunanya dan masih perlu dimatangkan lagi rancangan sistemnya. Implementasi sistem dijadwalkan pada tahap (B4–B6), yang mencakup pembangunan modul utama perangkat lunak audit energi sesuai desain yang telah ditetapkan. Pengujian dan evaluasi sistem akan dilakukan pada tahap (B7), yang bertujuan untuk memverifikasi dan memvalidasi fungsi serta kinerja sistem. Tahap akhir (B8) difokuskan pada penyempurnaan sistem dan penyusunan dokumentasi akhir. Pembagian waktu ini dirancang agar setiap tahap dapat dilaksanakan secara terfokus dan berurutan, sehingga mendukung pencapaian tujuan penelitian secara optimal.

\section{Analisis Risiko dan Mitigasi}
Pengembangan perangkat lunak audit energi tidak terlepas dari berbagai potensi risiko yang dapat memengaruhi kelancaran pelaksanaan dan kualitas hasil penelitian. Risiko tersebut dapat muncul pada aspek data, teknis, waktu, maupun ruang lingkup pengembangan sistem. Oleh karena itu, diperlukan identifikasi risiko sejak tahap perencanaan agar langkah mitigasi yang tepat dapat disiapkan. Ringkasan risiko utama yang teridentifikasi beserta dampak dan strategi mitigasinya disajikan pada Tabel~\ref{tbl:risks}.

\input table/tabelrisks.tex

Berdasarkan Tabel~\ref{tbl:risks}, risiko yang muncul sebagian besar berkaitan dengan keterbatasan data audit, perubahan kebutuhan pengguna, kendala teknis selama implementasi, serta keterbatasan waktu pengembangan. Strategi mitigasi yang dirancang berfokus pada pembatasan ruang lingkup penelitian, pemanfaatan data studi kasus yang tersedia, penyederhanaan fitur nonkritis, serta pengelolaan waktu pengembangan secara bertahap. Pendekatan ini diharapkan dapat meminimalkan dampak risiko terhadap proses pengembangan dan memastikan bahwa tujuan penelitian tetap dapat dicapai sesuai dengan rencana yang telah ditetapkan.